\section{Screening at the kinematic edge}\label{sec:15.2}

A particle at the threshold mass $m_{\mathrm{min}} = \SI{0.10}{keV}$ of \cref{we:7-al} must move at $v_{\mathrm{max}}$ and give up its whole kinetic energy, which takes the momentum transfer $q = m_{\mathrm{min}}v_{\mathrm{max}} = 2\omega/v_{\mathrm{max}}$, twice the edge value. Heavier particles can deposit the same energy with less momentum: from $v_{\mathrm{min}}(q,\omega) = v_{\mathrm{max}}$ the smallest momentum transfer at $\omega = 2\Delta_{\mathrm{BCS}}$ is 1.17 times the edge value for $2m_{\mathrm{min}}$ and only 1.03 times it at \SI{1}{keV}. For every particle a few times heavier than $m_{\mathrm{min}}$, the scatters just above $2\Delta_{\mathrm{BCS}}$ therefore lie close to the kinematic edge.

\subsection{The normal dielectric function on the edge}\label{sec:15.2-resum}

At $q \ll k_{\mathrm{F}}$ the Lindhard function \cref{eq:4.5-lindhard} depends on $\omega$ and $q$ only through the reduced frequency $\mathfrak{w} = \omega/qv_{\mathrm{F}}$, the common limit of the arguments $\omega/qv_{\mathrm{F}} \pm q/2k_{\mathrm{F}}$ as $q/2k_{\mathrm{F}} \to 0$. To first order in $q/k_{\mathrm{F}}$ the bracket of \cref{eq:4.5-lindhard} is the derivative $\partial_{\mathfrak{w}}[(1 - \mathfrak{w}^2)\ell(\mathfrak{w})] = 2 - 2\mathfrak{w}\,\ell(\mathfrak{w})$, with $\ell(\mathfrak{w}) = \ln[(\mathfrak{w} + 1)/(\mathfrak{w} - 1)]$ and $\ell' = 2/(1 - \mathfrak{w}^2)$. With $\varepsilon_{\mathrm{Lin}} = 1 - v_{\mathrm{C}}R^0$ and $k_{\mathrm{TF}}^2 = 2e^2\mathcal{N}_{\mathrm{F}}$, this gives $R^0 = -2\mathcal{N}_{\mathrm{F}}\mathfrak{f}_{\mathrm{L}}(\mathfrak{w})$ and, on the real axis,
\begin{equation}
  \begin{gathered}
  \varepsilon_{\mathrm{Lin}}(q,\omega) = 1 + \frac{k_{\mathrm{TF}}^2}{q^2}\Bigl[\Real\mathfrak{f}_{\mathrm{L}}(\mathfrak{w}) + \ii\,\frac{\pi\mathfrak{w}}{2}\Bigr] , \qquad
  \mathfrak{f}_{\mathrm{L}}(\mathfrak{w}) = 1 - \frac{\mathfrak{w}}{2}\,\ell(\mathfrak{w}) , \\
  \Real\mathfrak{f}_{\mathrm{L}} = 1 - \frac{\mathfrak{w}}{2}\ln\frac{1 + \mathfrak{w}}{1 - \mathfrak{w}} \qquad (\mathfrak{w} < 1) .
  \end{gathered}
  \label{eq:15.2-Lin}
\end{equation}
The function $\mathfrak{f}_{\mathrm{L}}$ is analytic in the upper half plane of $\mathfrak{w}$; at real $\mathfrak{w} < 1$ the retarded branch of $\ell$ gives the imaginary part $\pi\mathfrak{w}/2$. It has $\mathfrak{f}_{\mathrm{L}}(0) = 1$, the static limit, and $\mathfrak{f}_{\mathrm{L}} \to -1/3\mathfrak{w}^2$ for $|\mathfrak{w}| \gg 1$, the free-particle tail of \cref{eq:4.5-moments}. The kinematics of \cref{sec:1.3} require $q \ge \omega/v_{\mathrm{max}}$, so $\mathfrak{w} \le v_{\mathrm{max}}/v_{\mathrm{F}}$, with equality on the edge. In aluminum $v_{\mathrm{max}}/v_{\mathrm{F}} = 0.392$, and on the edge $\Real\mathfrak{f}_{\mathrm{L}} = 0.837$ at every $\omega$. The real part of the normal dielectric function there is \SI{16}{\percent} below the Thomas--Fermi value $k_{\mathrm{TF}}^2/q^2$, and the ratio of imaginary to real part,
\begin{equation*}
  \mathfrak{r}_\varepsilon = \frac{\Imag\varepsilon_{\mathrm{Lin}}}{\Real\varepsilon_{\mathrm{Lin}}} = \frac{\pi\mathfrak{w}}{2\Real\mathfrak{f}_{\mathrm{L}}(\mathfrak{w})} = 0.736 ,
\end{equation*}
is not small. With the ``1'' of $\varepsilon$ dropped, which changes nothing at these momenta (\cref{sec:15.1-rate}), every loss-function ratio below depends on $y = \omega/\Delta_{\mathrm{BCS}}$ and $qv_{\mathrm{F}}/\Delta_{\mathrm{BCS}}$ alone, and on the edge on $y$ and $v_{\mathrm{max}}/v_{\mathrm{F}}$. The aluminum numbers therefore apply to any clean weak-coupling superconductor with the same velocity ratio.

\begin{table}[tbp]
  \caption{Loss function of aluminum on the kinematic edge $q = \omega/v_{\mathrm{max}}$ ($\omega/qv_{\mathrm{F}} = v_{\mathrm{max}}/v_{\mathrm{F}} = 0.392$), divided by $\mathfrak{g}_{\mathrm{I}}L^{\mathrm{nor}}$, the value with normal-state screening, against $y = \omega/\Delta_{\mathrm{BCS}}$. Columns: case-I ratio $\mathfrak{g}_{\mathrm{I}}$ of \cref{eq:15.1-gI}; vertex fraction $\mathfrak{v}$ of \cref{sec:15.1-bare}; conserving response \cref{eq:15.1-Rcons}, clean, and with Born scattering by impurities at the mean free paths $\mathfrak{l}_e = \SI{1}{\micro m}$ and \SI{30}{nm} (\cref{sec:15.3-sc}), each divided by $\mathfrak{g}_{\mathrm{I}}$ times the normal state at the same scattering rate. All at $T = 0$, constant density of states, $q \ll k_{\mathrm{F}}$.}
  \label{tab:15-edge}
  \centering\small\setlength{\tabcolsep}{4pt}
  \begin{tabular}{@{}cccccc@{}}
    \toprule
    & & & \multicolumn{3}{c}{conserving} \\
    \cmidrule(l){4-6}
    $y$ & $\mathfrak{g}_{\mathrm{I}}$ & $\mathfrak{v}$ & clean & \SI{1}{\micro m} & \SI{30}{nm} \\
    \midrule
    2.05 & 1.552 & 0.22 & 0.681 & 0.365 & 0.045 \\
    2.2 & 1.500 & 0.44 & 0.867 & 0.600 & 0.187 \\
    2.5 & 1.418 & 0.35 & 0.999 & 0.854 & 0.569 \\
    3 & 1.322 & 0.21 & 1.069 & 1.021 & 1.008 \\
    5 & 1.151 & 0.052 & 1.067 & 1.057 & 1.091 \\
    10 & 1.051 & 0.004 & 1.023 & 1.020 & 1.021 \\
    \bottomrule
  \end{tabular}
\end{table}

\subsection{The conserving response on the edge}\label{sec:15.2-cons}

Let a response have $\Imag\varepsilon = \mathfrak{g}\,\Imag\varepsilon_{\mathrm{Lin}}$ and $\Real\varepsilon = (1 + \delta)\Real\varepsilon_{\mathrm{Lin}}$ at some $(q,\omega)$, with $\mathfrak{g}$ and $\delta$ real. Dividing $L = \Imag\varepsilon/[(\Real\varepsilon)^2 + (\Imag\varepsilon)^2]$ by the same expression for $\varepsilon_{\mathrm{Lin}}$, and then numerator and denominator by $(\Real\varepsilon_{\mathrm{Lin}})^2$,
\begin{equation}
  \frac{L}{L^{\mathrm{nor}}} = \mathfrak{g}\,\frac{1 + \mathfrak{r}_\varepsilon^2}{(1 + \delta)^2 + \mathfrak{g}^2\mathfrak{r}_\varepsilon^2} ,
  \label{eq:D.39-resum}
\end{equation}
with no approximation. The conserving response of \cref{eq:15.1-Rcons} enters with
\begin{equation}
  \mathfrak{g} = \mathfrak{g}_{\mathrm{I}}(1 + \mathfrak{v}) , \qquad
  \delta = \frac{\Real(\varepsilon^{\mathrm{cons}} - \varepsilon_{\mathrm{Lin}})}{\Real\varepsilon_{\mathrm{Lin}}} ,
  \label{eq:15.2-cons}
\end{equation}
with $\mathfrak{v}$ the vertex fraction of \cref{sec:15.1-bare}; both follow from \cref{eq:15.1-Rcons} by one quadrature over $\mu$. This separates what the pair-field ladder does to the two parts of $\varepsilon$, and on the edge it does more to the real part. At $y = 3$ the ladder raises the absorptive part from $1.322$ to $1.322 \times 1.214 = 1.605$ times normal, while it moves the real part from 1.33 times normal without the ladder to 0.60 times normal. The modulus of $\varepsilon$ ends close to the normal one, and the loss function close to $\mathfrak{g}_{\mathrm{I}}L^{\mathrm{nor}}$: $L^{\mathrm{cons}} = 1.07\,\mathfrak{g}_{\mathrm{I}}L^{\mathrm{nor}} = 1.41\,L^{\mathrm{nor}}$. Nearer the threshold the shift overshoots, and at $y = 2.2$ the real part of $\varepsilon^{\mathrm{cons}}$ is negative, $1 + \delta = -0.18$: the condensate responds inductively there. The superconductor nonetheless screens about as strongly as the normal metal, with $|\varepsilon^{\mathrm{cons}}|/|\varepsilon_{\mathrm{Lin}}| = 1.29$ at $y = 2.2$, 1.07 at 3 and 0.99 at 5. Along the whole edge the conserving value lies within \SI{10}{\percent} of $\mathfrak{g}_{\mathrm{I}}L^{\mathrm{nor}}$ for $y \ge 2.25$ and above it from $y = 2.5$ on, by at most \SI{8.4}{\percent} at $y = 3.6$; $L^{\mathrm{cons}}/L^{\mathrm{nor}}$ peaks at 1.43 at $y = 2.70$. These are clean values; at $y \le 2.5$ they hold only for mean free paths $\mathfrak{l}_e \gtrsim \SI{10}{\micro m}$ (\cref{sec:15.3-sc}).

Near the threshold, $y \lesssim 2.3$, the ladder is large because the Anderson--Bogoliubov pole of \cref{eq:15.1-Rcons} lies just below the gap, at $1.927$, $1.966$ and $1.992\,\Delta_{\mathrm{BCS}}$ for the edge momenta of $y = 2.2$, 2.5 and 3 (\cref{sec:15.1-closed}). At $y \ge 3$ the shift comes from the pair continuum of the ladder instead. The relative size $|R^{\mathrm{cons}} - R^0_{\boldsymbol{\eta}\boldsymbol{\eta}}|/|R^0_{\boldsymbol{\eta}\boldsymbol{\eta}}|$ is $27(\Delta_{\mathrm{BCS}}/qv_{\mathrm{F}})^2$ on the edge at $y = 3$, and its coefficient falls only slowly with $q$, to 18 at $qv_{\mathrm{F}} = 100\,\Delta_{\mathrm{BCS}}$ and 11 at $10^3\,\Delta_{\mathrm{BCS}}$.

\begin{workedexample}{The conserving loss function of aluminum at $y = 3$}\label{we:15-edge}
Take $\omega = 3\Delta_{\mathrm{BCS}} = \SI{0.54}{meV}$ with the bulk gap $\Delta_{\mathrm{BCS}} = \SI{0.18}{meV}$. The smallest momentum transfer is $q = \omega/v_{\mathrm{max}} = \SI{0.54e-3}{eV}/2.650\times10^{-3} = \SI{0.204}{eV}$, so $1/q = \SI{0.97}{\micro m}$. With $v_{\mathrm{F}} = 6.75\times10^{-3}$, $qv_{\mathrm{F}} = \SI{1.38}{meV} = 7.65\,\Delta_{\mathrm{BCS}}$, so $\mathfrak{w} = 0.392$ and $\Delta_{\mathrm{BCS}}/qv_{\mathrm{F}} = 0.131$.

\emph{Normal metal.} \Cref{eq:15.2-Lin} gives $\Real\mathfrak{f}_{\mathrm{L}} = 1 - 0.196\ln(1.392/0.608) = 0.837$ and $\mathfrak{r}_\varepsilon = (\pi/2)(0.392)/0.837 = 0.736$.

\emph{Conserving response.} The case-I factor is $\mathfrak{g}_{\mathrm{I}}(3) = 1.322$, and \cref{eq:15.2-cons} has $\mathfrak{g} = 1.322\times1.214 = 1.605$ and $1 + \delta = 0.596$. Then $|\varepsilon^{\mathrm{cons}}|/|\varepsilon_{\mathrm{Lin}}| = [(0.596^2 + 1.605^2\times0.736^2)/(1 + 0.736^2)]^{1/2} = 1.065$, and $L^{\mathrm{cons}} = 1.605\,L^{\mathrm{nor}}/1.065^2 = 1.414\,L^{\mathrm{nor}} = 1.069\,\mathfrak{g}_{\mathrm{I}}L^{\mathrm{nor}}$. In a \SI{30}{nm} film the ratio is 1.008 (\cref{tab:15-edge}).

The conserving response predicts a loss function \SI{41}{\percent} above the normal metal at this energy and momentum, almost all of it the case-I factor, and the mean free path changes that by less than \SI{6}{\percent}.
\end{workedexample}

\subsection{Corrections and nuisances}\label{sec:15.2-nuisance}

\Cref{eq:15.2-cons} inherits the hypotheses of \cref{eq:15.1-Rcons}, and \cref{sec:15.1-rate,sec:15.1-closed} bound the terms it drops: the band curvature, the ``1'' and the core electrons in $\varepsilon$, and the particle--hole asymmetry that couples the amplitude channel, each below $10^{-4}$ at these momenta. Thermal quasiparticles add the scattering channel of \cref{sec:7.2} with weight $\ee^{-\Delta_{\mathrm{BCS}}/T}$, $\ee^{-209}$ at \SI{10}{mK}. Within the model every correction to the absorptive part of order $(\Delta_{\mathrm{BCS}}/qv_{\mathrm{F}})^2$ is in the ladder, since the bare absorptive part has none, and \cref{eq:15.1-Rcons} contains the ladder to all orders. Outside the model lie disorder (\cref{sec:15.3-sc}) and the frequency dependence of the gap in strong coupling, which is small in aluminum, where $2\Delta_{\mathrm{BCS}}/T_c = 3.54$ against the weak-coupling 3.53 (\cref{sec:7.2}), and is not sized here. So is the film geometry, once the length $1/q$, about a micrometer on the edge, exceeds the film thickness (\cref{sec:15.3-sheet}).

The superconducting part of a rate carries three nuisances of the kinds of \cref{sec:8.8}, which the chapter summary collects. The first is the gap, a threshold nuisance \cref{eq:8.8-shift}: $\delta E = 2\,\delta\Delta_{\mathrm{BCS}}$ moves the edge, and at fixed $\omega$ it moves $y$. Its size is the uncertainty of the gap of the sensor in use, the bulk \SI{0.18}{meV} or the measured film gap of \cref{sec:15.1-al}. The second is the relaxation rate $\gamma_{\mathrm{rel}}$ of the film, a tilt that acts only near the edge. At $y \ge 3$ it moves the loss function by less than \SI{6}{\percent} relative to $\mathfrak{g}_{\mathrm{I}}L^{\mathrm{nor}}$, and at $y = 2.2$ by a factor of 4.6 between the clean limit and $\mathfrak{l}_e = \SI{30}{nm}$ (\cref{sec:15.3-sc}). The third is the edge remainder, the difference between the loss function a rate calculation uses and the conserving one. With $\mathfrak{g}_{\mathrm{I}}L^{\mathrm{nor}}$ as the central value, the remainder at $y \ge 3$ has no fixed sign. In the clean limit it runs from $-\SI{2.0}{\percent}$ to $+\SI{8.5}{\percent}$ over all momentum transfers that halo dark matter reaches, the lower end at about four times the edge momentum of $y = 3$. With impurities it reaches $+\SI{9.1}{\percent}$ on the edge, and off the edge it has not been computed. It is therefore a two-sided normalization nuisance, $-\SI{2}{\percent}$ to $+\SI{9}{\percent}$, with the standard-normal constraint of \cref{eq:8.8-M}.

Near the threshold the conserving loss function depends strongly on the mean free path, and a rate calculation excludes those energies from its signal region. The criterion we adopt is that the edge value of $L^{\mathrm{cons}}/\mathfrak{g}_{\mathrm{I}}L^{\mathrm{nor}}$ vary by at most \SI{10}{\percent} over the mean free paths from \SI{30}{nm} to \SI{100}{\micro m}. It holds at $y = 3$, where the values agree to \SI{6}{\percent}, and fails at $y = 2.5$, where they span a factor of 1.75. The boundary lies between the two, and the energies $y < 3$ are excluded. In that window the remainder is not bounded for a film: at $\mathfrak{l}_e = \SI{30}{nm}$ the conserving value at $y = 2.2$ is 0.19. In a crystal with $\mathfrak{l}_e \ge \SI{10}{\micro m}$ the clean column of \cref{tab:15-edge} holds to \SI{4}{\percent} down to $y = 2.2$, and the excluded energies can be restored with it. The deviations from $\mathfrak{g}_{\mathrm{I}}L^{\mathrm{nor}}$ are largest on the edge, so the exclusion covers the momenta above it as well.
